\documentclass[a4paper,11pt]{article}
\usepackage[utf8]{inputenc}
\usepackage[T1]{fontenc}
\usepackage[french]{babel}
\usepackage[left=2cm,right=2cm,top=2cm,bottom=2cm]{geometry}
\usepackage{xcolor}
\usepackage{hyperref}
\usepackage{fontawesome5}
\usepackage{parskip}

% --- Couleurs EDF ---
\definecolor{primary}{RGB}{0, 85, 164}
\hypersetup{colorlinks=true, urlcolor=primary}

\begin{document}

% --- EXPÉDITEUR ---
\noindent
\begin{minipage}[t]{0.5\textwidth}
    \textbf{Loïc Aron MBASSI EWOLO}\\
    Élève-Ingénieur Bac+5 -- Double diplôme ENSEA\\[3pt]
    \faMapMarker\ Cergy, France\\
    \faPhone\ (+33) 7 59 52 33 98\\
    \faEnvelope\ \href{mailto:loicaron.mbassiewolo@ensea.fr}{loicaron.mbassiewolo@ensea.fr}
\end{minipage}
\hfill
% --- DATE ---
\begin{minipage}[t]{0.4\textwidth}
    \raggedleft
    Cergy, le 27 février 2026
\end{minipage}

\vspace{1.2cm}

% --- DESTINATAIRE ---
\hfill
\begin{minipage}[t]{0.5\textwidth}
    \raggedleft
    \textbf{EDF -- DOAAT}\\
    Service Recrutement / Équipe SoDATA\\
    Saint-Denis, France
\end{minipage}

\vspace{1cm}

\noindent\textbf{Objet :} Candidature en alternance -- Data Scientist IA (12 mois dès septembre 2026)

\vspace{0.8cm}

Madame, Monsieur,

Concevoir un système multi-agents capable d'orchestrer des LLM, d'interroger un datalake via MCP et de produire des visualisations automatiques : c'est le type de défi technique sur lequel je travaille depuis plus d'un an. En 2024, dans le cadre d'un projet de recherche, j'ai conçu un système de 4 agents IA collaboratifs utilisant Gemini et RAG pour l'enrichissement contextuel, orchestrés via Google ADK, avec intégration d'APIs externes (météo, marchés) et déploiement FastAPI/Docker. Ce projet recouvre directement les enjeux de l'assistant multi-agents que l'équipe SoDATA souhaite développer.

En parallèle de cette expérience en IA agentique, j'ai travaillé comme assistant de recherche au MILA (Institut Québécois d'IA) sur l'interprétabilité des réseaux de neurones, ce qui m'a donné une rigueur dans l'analyse des comportements de modèles et la reproductibilité des expériences. Mon travail sur le fine-tuning de Transformers (BERT, GPT-2) lors de la compétition MLPC et mes certifications en Machine Learning (Stanford/DeepLearning.AI) complètent ma maîtrise de l'écosystème LLM que vous utilisez (LangChain, OpenAI, prompting, RAG).

Sur le volet data et visualisation, j'ai remporté le 1\textsuperscript{er} prix IndabaX Africa 2025 en construisant un pipeline complet : nettoyage de données médicales, déploiement d'API Flask et dashboard Streamlit interactif. Cette capacité à transformer des données brutes en visualisations exploitables par des non-techniciens correspond aux besoins de génération automatique de graphes et rapports de votre assistant. Mon profil mathématique (algèbre linéaire, probabilités, statistiques) me permet également de challenger les résultats produits par les agents en termes de précision et de cohérence.

Je suis en double diplôme ENSEA / Polytechnique Yaoundé, avec une spécialisation en IA et traitement du signal. Je termine actuellement un stage de recherche en neurosciences computationnelles à Florence (Italie) jusqu'en août 2026, où j'exploite des publications scientifiques en anglais et propose des axes de recherche, et je suis disponible dès septembre 2026 pour une alternance de 12 mois. Le rythme proposé par l'ENSEA est de 3 semaines en entreprise / 1 semaine en cours. Contribuer à l'adoption du MCP comme standard d'interopérabilité au sein du Groupe EDF, tout en développant mes compétences en IA agentique dans un contexte industriel à fort impact, correspond pleinement à mon projet professionnel.

Je reste à votre disposition pour un entretien afin de discuter de ma candidature.

Je vous prie d'agréer, Madame, Monsieur, l'expression de mes salutations distinguées.

\vspace{0.8cm}

\textbf{Loïc Aron MBASSI EWOLO}\\
\textit{Élève-Ingénieur ENSEA / Polytechnique Yaoundé}

\end{document}
